\begin{filecontents*}{\jobname.bib}
@article{ErdosHajnal1989,
  author  = {Erd\H{o}s, Paul and Hajnal, Andr\'as},
  title   = {Ramsey-type theorems},
  journal = {Discrete Applied Mathematics},
  volume  = {25},
  number  = {1--2},
  pages   = {37--52},
  year    = {1989}
}

@article{ChudnovskySurvey2014,
  author  = {Chudnovsky, Maria},
  title   = {The {Erd\H{o}s-Hajnal} conjecture: A survey},
  journal = {Journal of Graph Theory},
  volume  = {75},
  number  = {2},
  pages   = {178--190},
  year    = {2014},
  doi     = {10.1002/jgt.21730}
}

@article{AlonPachSolymosi2001,
  author  = {Alon, Noga and Pach, J\'anos and Solymosi, J\'ozsef},
  title   = {Ramsey-type theorems with forbidden subgraphs},
  journal = {Combinatorica},
  volume  = {21},
  number  = {2},
  pages   = {155--170},
  year    = {2001},
  doi     = {10.1007/s004930100016}
}

@article{BucicNguyenScottSeymour2024DensityI,
  author  = {Buci\'c, Matija and Nguyen, Tung and Scott, Alex and Seymour, Paul},
  title   = {Induced subgraph density. {I}. A loglog step towards {Erd\H{o}s-Hajnal}},
  journal = {International Mathematics Research Notices},
  volume  = {2024},
  number  = {12},
  pages   = {9991--10004},
  year    = {2024},
  doi     = {10.1093/imrn/rnad301}
}

@misc{BucicFoxPham2024,
  author       = {Buci\'c, Matija and Fox, Jacob and Pham, Huy Tuan},
  title        = {Equivalence between {Erd\H{o}s-Hajnal} and polynomial {R\"odl} and {Nikiforov} conjectures},
  year         = {2024},
  eprint       = {2403.08303},
  archivePrefix= {arXiv},
  primaryClass = {math.CO}
}

@article{ChudnovskySafra2008Bull,
  author  = {Chudnovsky, Maria and Safra, Shmuel},
  title   = {The {Erd\H{o}s-Hajnal} conjecture for bull-free graphs},
  journal = {Journal of Combinatorial Theory, Series B},
  volume  = {98},
  number  = {6},
  pages   = {1301--1310},
  year    = {2008},
  doi     = {10.1016/j.jctb.2008.02.005}
}

@article{ChudnovskyScottSeymourSpirkl2023C5,
  author  = {Chudnovsky, Maria and Scott, Alex and Seymour, Paul and Spirkl, Sophie},
  title   = {{Erd\H{o}s-Hajnal} for graphs with no 5-hole},
  journal = {Proceedings of the London Mathematical Society},
  volume  = {126},
  number  = {3},
  pages   = {997--1014},
  year    = {2023},
  doi     = {10.1112/plms.12504}
}

@article{NguyenScottSeymour2026P5,
  author  = {Nguyen, Tung and Scott, Alex and Seymour, Paul},
  title   = {Induced subgraph density. {VII}. The five-vertex path},
  journal = {Proceedings of the London Mathematical Society},
  volume  = {132},
  number  = {3},
  pages   = {e70133},
  year    = {2026},
  doi     = {10.1112/plms.70133}
}

@misc{NguyenScottSeymour2023DensityIV,
  author       = {Nguyen, Tung and Scott, Alex and Seymour, Paul},
  title        = {Induced subgraphs density. {IV}. New graphs with the {Erd\H{o}s-Hajnal} property},
  year         = {2023},
  eprint       = {2307.06455},
  archivePrefix= {arXiv},
  primaryClass = {math.CO}
}

@misc{NguyenScottSeymour2023DensityV,
  author       = {Nguyen, Tung and Scott, Alex and Seymour, Paul},
  title        = {Induced subgraph density. {V}. All paths approach {Erd\H{o}s-Hajnal}},
  year         = {2023},
  eprint       = {2307.15032},
  archivePrefix= {arXiv},
  primaryClass = {math.CO}
}

@article{NguyenScottSeymour2025DensityVI,
  author  = {Nguyen, Tung and Scott, Alex and Seymour, Paul},
  title   = {Induced subgraph density. {VI}. Bounded {VC}-dimension},
  journal = {Advances in Mathematics},
  volume  = {482},
  pages   = {110601},
  year    = {2025},
  doi     = {10.1016/j.aim.2025.110601}
}

@article{ChudnovskyKimOumSeymour2016,
  author  = {Chudnovsky, Maria and Kim, Ringi and Oum, Sang-il and Seymour, Paul},
  title   = {Unavoidable induced subgraphs in large graphs with no homogeneous sets},
  journal = {Journal of Combinatorial Theory, Series B},
  volume  = {118},
  pages   = {1--12},
  year    = {2016},
  doi     = {10.1016/j.jctb.2016.01.008}
}

@article{MalliarisTerry2018,
  author  = {Malliaris, Maryanthe and Terry, Caroline},
  title   = {On unavoidable induced subgraphs in large prime graphs},
  journal = {Journal of Graph Theory},
  volume  = {88},
  number  = {2},
  pages   = {255--270},
  year    = {2018},
  doi     = {10.1002/jgt.22209}
}

@article{Rodl1986,
  author  = {R\"odl, Vojt\v{e}ch},
  title   = {On universality of graphs with uniformly distributed edges},
  journal = {Discrete Mathematics},
  volume  = {59},
  pages   = {125--134},
  year    = {1986},
  doi     = {10.1016/0012-365X(86)90076-2}
}

@article{FoxSudakov2008,
  author  = {Fox, Jacob and Sudakov, Benny},
  title   = {Induced {Ramsey}-type theorems},
  journal = {Advances in Mathematics},
  volume  = {219},
  number  = {6},
  pages   = {1771--1800},
  year    = {2008},
  doi     = {10.1016/j.aim.2008.07.009}
}

@article{BousquetLagoutteThomasse2015,
  author  = {Bousquet, Nicolas and Lagoutte, Aur\'elie and Thomass\'e, St\'ephan},
  title   = {The {Erd\H{o}s-Hajnal} conjecture for paths and antipaths},
  journal = {Journal of Combinatorial Theory, Series B},
  volume  = {113},
  pages   = {261--264},
  year    = {2015},
  doi     = {10.1016/j.jctb.2015.01.001}
}

@article{AlecuAtminasLozinZamaraev2021,
  author  = {Alecu, Bogdan and Atminas, Aistis and Lozin, Vadim V. and Zamaraev, Viktor},
  title   = {Graph classes with linear {Ramsey} numbers},
  journal = {Discrete Mathematics},
  volume  = {344},
  number  = {4},
  pages   = {112307},
  year    = {2021},
  doi     = {10.1016/j.disc.2021.112307}
}

@article{ErdosSzekeres1935,
  author  = {Erd\H{o}s, Paul and Szekeres, George},
  title   = {A combinatorial problem in geometry},
  journal = {Compositio Mathematica},
  volume  = {2},
  pages   = {463--470},
  year    = {1935}
}

@article{Sauer1972,
  author  = {Sauer, Norbert},
  title   = {On the density of families of sets},
  journal = {Journal of Combinatorial Theory, Series A},
  volume  = {13},
  number  = {1},
  pages   = {145--147},
  year    = {1972},
  doi     = {10.1016/0097-3165(72)90019-2}
}

@article{Shelah1972,
  author  = {Shelah, Saharon},
  title   = {A combinatorial problem; stability and order for models and theories in infinitary languages},
  journal = {Pacific Journal of Mathematics},
  volume  = {41},
  number  = {1},
  pages   = {247--261},
  year    = {1972},
  doi     = {10.2140/pjm.1972.41.247}
}
\end{filecontents*}

\documentclass[11pt]{amsart}
\usepackage[T1]{fontenc}
\usepackage[utf8]{inputenc}
\usepackage{lmodern}
\usepackage{microtype}
\usepackage[margin=0.82in]{geometry}
\usepackage{amsmath,amssymb,amsthm,mathtools}
\usepackage{enumitem}
\usepackage[hidelinks]{hyperref}
\usepackage[nameinlink,noabbrev]{cleveref}
\usepackage[backend=biber,style=alphabetic,maxbibnames=99,giveninits=true]{biblatex}
\addbibresource{\jobname.bib}
\setlength{\emergencystretch}{4em}
\setlength{\parskip}{0.08em}
\linespread{1.02}
\setlist{leftmargin=*,itemsep=0.22em,topsep=0.35em}
\numberwithin{equation}{section}

\newtheorem{theorem}{Theorem}[section]
\newtheorem{proposition}[theorem]{Proposition}
\newtheorem{lemma}[theorem]{Lemma}
\newtheorem{corollary}[theorem]{Corollary}
\theoremstyle{definition}
\newtheorem{definition}[theorem]{Definition}
\newtheorem{question}[theorem]{Question}
\theoremstyle{remark}
\newtheorem{remark}[theorem]{Remark}

\newcommand{\F}{\mathbb F}
\newcommand{\VCdim}{\operatorname{VCdim}}
\DeclareMathOperator{\im}{im}

\title[Conservative obstructions toward EH]{Conservative obstructions and local checks around several Erd\H{o}s-Hajnal routes}
\author{Samuel Mausberg}
\date{April 2026, forum version}

\begin{document}
\begin{abstract}
This note is a trimmed and conservative version of working notes around the Erd\H{o}s-Hajnal conjecture. It does not claim to prove the conjecture, and it does not claim priority for any standard-looking lemma. The aim is to isolate a few locally checkable obstructions to tempting reductions: a singleton-cut obstruction to bounded affine-cell occupancy in level-two shadows, exact path-profile and first-splitter obstructions outside a witness path, a semantic gap after ordered cutpoint canonization, a correction to a two-anchor blockade notion, and a local rooted $P_4$ crossing-trace obstruction. The emphasis is on statements that are short enough to be checked or refuted by readers, with folklore-sensitive framing and explicit open questions.
\end{abstract}

\maketitle

\section{Scope and literature placement}

The Erd\H{o}s-Hajnal conjecture asserts that for every graph $H$ there is $c(H)>0$ such that every $H$-free graph $G$ has
\[
  \max\{\omega(G),\alpha(G)\}\ge |G|^{c(H)}.
\]
The conjecture is open in general. Background and early reductions are surveyed in \cite{ChudnovskySurvey2014}. The substitution theorem of Alon, Pach, and Solymosi reduces the conjecture to prime forbidden graphs \cite{AlonPachSolymosi2001}. Known five-vertex cases now include the bull, $C_5$, and $P_5$ \cite{ChudnovskySafra2008Bull,ChudnovskyScottSeymourSpirkl2023C5,NguyenScottSeymour2026P5}. The best general lower bound was improved by Buci\'c, Nguyen, Scott, and Seymour \cite{BucicNguyenScottSeymour2024DensityI}. Buci\'c, Fox, and Pham proved equivalence between the Erd\H{o}s-Hajnal conjecture and polynomial forms of R\"odl and Nikiforov type conjectures \cite{BucicFoxPham2024}, putting polynomial pure-pair and restricted-induced-subgraph statements near the center of the current picture. Large prime graphs are constrained by the unavoidable-template theorem of Chudnovsky, Kim, Oum, and Seymour and the model-theoretic treatment of Malliaris and Terry \cite{ChudnovskyKimOumSeymour2016,MalliarisTerry2018}.

This note is best read in that context. It is not a new global strategy. It records several places where natural local strategies appear to need an additional $H$-dependent input. The current positive literature already contains much stronger machinery than the elementary pieces below: iterative sparsification and buildable graphs \cite{NguyenScottSeymour2023DensityIV}, near-Erd\H{o}s-Hajnal bounds for all paths \cite{NguyenScottSeymour2023DensityV}, bounded VC-dimension \cite{NguyenScottSeymour2025DensityVI}, and path-antipath results \cite{BousquetLagoutteThomasse2015}. The point here is narrower: when a proposed route uses affine shadows, path profiles, cutpoint canonization, blockades, or rooted $P_4$ two-anchor extraction, the following examples indicate what must still be proved.

\paragraph{Provenance and novelty status.} This note was developed from AI-assisted exploratory conversations and subsequent human editing. It is meant to be falsifiable and locally checkable. Several ingredients below are elementary and may be folklore, especially the affine/co-point classification, the cutpoint canonization lemma, and the rectangle shattering lemma. They are included because they clarify exact failure modes, not because they should be read as priority claims. The statements most worth external scrutiny are \Cref{thm:singleton-cut}, \Cref{thm:first-splitter,prop:chain-quotients}, \Cref{prop:two-anchor-not-rich,prop:coordinate-synchronizes}, and \Cref{thm:H-specific-cubes}.

All graphs are finite and simple. For disjoint sets $A,B$, $A$ is complete to $B$ if every edge between $A$ and $B$ is present, and anticomplete to $B$ if no such edge is present. A graph is prime if it has no homogeneous set, where $X$ is homogeneous if $2\le |X|<|G|$ and every vertex outside $X$ is either complete or anticomplete to $X$.

\section{Level-two affine shadows and the occupancy obstruction}

This section records the cleanest obstruction in the note. The linear algebra inside an affine cell is useful but elementary. The nontrivial warning is that connected row-consistency and low ladder index do not force bounded affine-cell occupancy.

\subsection{Arity two: affine lines versus co-points}

Let $G$ be a graph, let $Y\subseteq V(G)$, and let $A\subseteq (V(G)\setminus Y)^{(2)}$ be a set of ordered pairs of distinct vertices. For $(u,v)\in A$ and $y\in Y$, define
\[
  \lambda((u,v),y)=(1_{uy\in E(G)},1_{vy\in E(G)})\in \F_2^2.
\]
A subset of $\F_2^2$ is affine if it is an affine subspace, and is a co-point if it is the complement of one point.

\begin{proposition}[Arity-two pullback and universality]\label{prop:arity-two}
Let $G,Y,A$ be as above.
\begin{enumerate}[label=(\alph*)]
\item If $\lambda(A\times Y)$ is contained in a proper affine subspace of $\F_2^2$, then $G$ contains disjoint sets $S,T$ with
\[
  |S|\ge |A|/|G|, \qquad |T|\ge |Y|/2,
\]
and $S$ complete or anticomplete to $T$.
\item If $C\subseteq \F_2^2$ is a co-point, then for arbitrary finite sets $I,J$ and arbitrary $\mu:I\times J\to C$, there is a graph with pairs $(u_i,v_i)$ and vertices $y_j$ such that $\lambda((u_i,v_i),y_j)=\mu(i,j)$ for all $i,j$.
\end{enumerate}
Consequently, for arbitrary pair-sets the nonempty proper images in $\F_2^2$ split exactly into affine pullback cases and universal co-point cases.
\end{proposition}

\begin{proof}
For (a), every proper affine subspace of $\F_2^2$ is contained in one of the affine lines $x_1=c$, $x_2=c$, or $x_1+x_2=c$. If $x_1=c$, the set of first coordinates of pairs in $A$ is complete or anticomplete to $Y$, and has size at least $|A|/|G|$; the case $x_2=c$ is symmetric. If $x_1+x_2=0$, partition $V(G)\setminus Y$ by exact neighborhoods in $Y$. Pairs in $A$ lie inside equal trace classes, so some class has size at least $|A|/|G|$. Inside that class, either the common neighborhood in $Y$ or its complement has size at least $|Y|/2$. If $x_1+x_2=1$, pairs join complementary trace classes, and the same counting gives one trace class of size at least $|A|/|G|$.

For (b), take vertices $u_i,v_i$ for $i\in I$ and $y_j$ for $j\in J$. Add $u_i y_j$ and $v_i y_j$ exactly according to the two coordinates of $\mu(i,j)$, and put no other relevant edges. This realizes the prescribed matrix.
\end{proof}

\subsection{Connected row-consistent shadows}

Let $P$ be a graph on vertex set $U$ and let $Y$ be finite. A labeling
\[
  \mu:E(P)\times Y\to \F_2^2
\]
is row-consistent if for each $y\in Y$ there is a function $a_y:U\to \F_2$ such that
\[
  \mu(uv,y)=(a_y(u),a_y(v))\qquad (uv\in E(P)).
\]
The parity trace is $\pi_y(uv)=a_y(u)+a_y(v)$. If $P$ is connected, then parity determines the exact row up to global complement: if $a,b:U\to \F_2$ have the same parity on every edge of $P$, then $a+b$ is constant on $U$.

For a finite set $Q$, an exact affine cell on $Q$ is a nonempty coset $X=\eta+W\subseteq \F_2^Q$. Its dimension is $\dim W$. For $E\subseteq \binom Q2$, define
\[
  \operatorname{tr}_E(x)=((x_p,x_q))_{\{p,q\}\in E},
  \qquad
  \partial_E(x)=(x_p+x_q)_{\{p,q\}\in E}.
\]
A binary matrix contains a ladder of size $k$ if it has rows $r_1,\dots,r_k$ and columns $c_1,\dots,c_k$ such that the $(i,j)$ entry is $1$ exactly when $i\le j$.

\begin{theorem}[Linear algebra inside an exact affine cell]\label{thm:affine-cell}
Let $Q$ be finite and let $X=\eta+W\subseteq \F_2^Q$ be an exact affine cell of dimension $r$.
\begin{enumerate}[label=(\roman*)]
\item For distinct $p,q\in Q$ and $(\alpha,\beta)\in \F_2^2$, the fiber
\[
  X^{\alpha,\beta}_{p,q}=X\cap\{x:x_p=\alpha, x_q=\beta\}
\]
is either empty or an exact affine cell. If nonempty and proper, its dimension is less than $r$.
\item If $r>0$ and $|Q|\ge 2$, then some such fiber is nonempty, proper, and has size at least $|X|/4$.
\item For every $E\subseteq \binom Q2$, $|\operatorname{tr}_E(X)|\le 2^r$.
\item The parity matrix $(\partial_E(x)_e)_{x\in X,e\in E}$ has ladder index at most $r+1$.
\end{enumerate}
\end{theorem}

\begin{proof}
Fibers are intersections of $X$ with two affine coordinate hyperplanes, proving (i). If $r>0$, some coordinate is nonconstant on $X$; pairing it with another coordinate gives a nonconstant map to $\F_2^2$, so one of at most four fibers is nonempty proper and has size at least $|X|/4$, proving (ii). The trace map is affine on the $r$-dimensional space $W$, giving (iii).

For (iv), suppose a ladder of size $k$ appears, with rows $x_1,\dots,x_k\in X$ and columns $e_1,\dots,e_k\in E$. Put $d_i=x_i+x_{i+1}\in W$ for $1\le i<k$. If $\ell_j$ is the linear part of the parity functional at $e_j$, then
\[
  \ell_j(d_i)=\partial_E(x_i)_{e_j}+\partial_E(x_{i+1})_{e_j}=1
  \quad\text{if and only if}\quad i=j.
\]
Thus $d_1,\dots,d_{k-1}$ are linearly independent in $W$, so $k-1\le r$.
\end{proof}

\begin{theorem}[Singleton-cut obstruction to affine occupancy]\label{thm:singleton-cut}
For every $m>4$ there is a connected row-consistent level-two parity shadow whose occupied rows lie in an affine cut space of dimension $m-1$, have ladder index at most $2$, and have harmless two-coordinate projections, but cannot be partitioned into $O(1)$ exact affine cells. More precisely, let $\Gamma=K_m$, let
\[
  \delta:\F_2^{[m]}\to \F_2^{E(K_m)},\qquad (\delta x)_{ij}=x_i+x_j,
\]
let $C_m=\im\delta$, let $c_i=\delta e_i$ for $i\in[m]$, let $c_0=0$, and let
\[
  R_m=\{c_0,c_1,\dots,c_m\}\subseteq C_m.
\]
Then:
\begin{enumerate}[label=(\roman*)]
\item $K_m$ is connected and every row of $R_m$ is cycle-consistent;
\item $\dim C_m=m-1$;
\item for any two edge coordinates $e,f\in E(K_m)$, the projection of $R_m$ to $\F_2^{\{e,f\}}$ is either all of $\F_2^2$ or the co-point $\{00,10,01\}$ up to coordinate relabeling;
\item every exact affine cell contained in $R_m$ has size at most $2$, so every exact affine-cell partition of $R_m$ has at least $\lceil(m+1)/2\rceil$ cells;
\item the parity matrix with rows $R_m$ and columns $E(K_m)$ has ladder index at most $2$.
\end{enumerate}
Thus connected row-consistency and bounded ladder index do not imply bounded exact affine-cell occupancy.
\end{theorem}

\begin{proof}
Each $c_i$ is a cut vector, so (i) holds. Since $K_m$ is connected, $\ker\delta=\langle \mathbf 1\rangle$, and $\dim C_m=m-1$. Equivalently, the only linear relation among $c_1,\dots,c_m$ is $c_1+\cdots+c_m=0$.

For (iii), fix two edge coordinates $e,f$. If $e$ and $f$ share an endpoint, the shared endpoint gives pattern $11$, the two nonshared endpoints give $10$ and $01$, and $c_0$ gives $00$. If $e$ and $f$ are disjoint, no singleton cut hits both edges, so the observed patterns are $00,10,01$.

For (iv), an affine plane over $\F_2$ consists of four points whose sum is zero. If four distinct points of $R_m$ include $c_0$, then the other three singleton cuts would sum to zero, impossible. If the four points are $c_i,c_j,c_k,c_\ell$ with distinct nonzero indices, their sum is $\delta 1_A$ for a four-element set $A$, which is nonzero because $m>4$. Hence $R_m$ contains no affine plane, so each affine cell contained in it has dimension at most one and size at most two.

For (v), each edge-column is supported on exactly two nonzero singleton cuts and not on $c_0$. In a ladder of size $k$, the last ladder column is $1$ on all $k$ selected rows, so $k\le 2$.
\end{proof}

\begin{remark}
The obstruction does not rule out an $H$-dependent affine route. It rules out the unconditional step ``connected row-consistent co-point data plus bounded ladder index gives bounded affine-cell occupancy.'' A valid route needs either a dense-cell extraction theorem using $H$-freeness, or an escape that turns sparse occupancy into a polynomial pure pair.
\end{remark}

\section{Path profiles and first splitters}

The results in this section are local checks around induced paths. They should not be confused with the global literature on paths: the Erd\H{o}s-Hajnal conjecture is known for $P_5$ and nearly true for all paths in the sense of \cite{NguyenScottSeymour2023DensityV,NguyenScottSeymour2026P5}. The point here is only to locate where simple witness-profile quotients stop.

\subsection{A terminal safe bag is not a proper-subgraph reduction}

\begin{proposition}[Terminal blow-up obstruction]\label{prop:terminal-blowup}
Fix $k\ge 4$. Let $C=c_1-\cdots-c_{k-1}$ be an induced $P_{k-1}$, let $W$ be any $P_k$-free graph disjoint from $C$, and form $G_k(W)$ by making every vertex of $W$ adjacent exactly to $c_{k-2}$ and $c_{k-1}$. Then $G_k(W)$ is $P_k$-free.
\end{proposition}

\begin{proof}
Let $P$ be an induced path in $G_k(W)$ that meets both $C$ and $W$. The only vertices of $C$ adjacent to $W$ are $c_{k-2}$ and $c_{k-1}$. The path $P$ cannot contain both $c_{k-2}$ and $c_{k-1}$ together with a vertex of $W$, since these three vertices form a triangle.

If $P$ contains $c_{k-1}$ but not $c_{k-2}$, then no earlier vertex of $C$ is connected to the $W$ part of $P$. Also $c_{k-1}$ is adjacent to every vertex of $W$, so an induced path through $c_{k-1}$ contains at most one vertex of $W$ beyond its neighbor. Thus this case has fewer than $k$ vertices. If $P$ contains $c_{k-2}$ but not $c_{k-1}$, then either the $C$ part is contained in $c_1-\cdots-c_{k-2}$ and at most one vertex of $W$ can be added when $c_{k-3}$ is present, giving at most $k-1$ vertices, or $c_{k-3}$ is absent and the mixed component has at most three vertices. If neither $c_{k-2}$ nor $c_{k-1}$ is present, the path cannot be mixed. Therefore no mixed induced $P_k$ exists. Since $W$ and $C$ themselves are $P_k$-free, the graph $G_k(W)$ is $P_k$-free.
\end{proof}

This corrects the common but false shortcut that a mixed path here cannot contain two vertices of $W$. It can contain two, but only in configurations too short to make a $P_k$.

\subsection{Exact safe profiles}

Fix $Q=q_1-\cdots-q_{k-1}$. For $S\subseteq V(Q)$, let $Q_S$ be obtained from $Q$ by adding one new vertex adjacent exactly to $S$.

\begin{lemma}\label{lem:one-vertex-extension}
The graph $Q_S$ is $P_k$-free if and only if $S\notin\{\{q_1\},\{q_{k-1}\}\}$.
\end{lemma}

\begin{proof}
If $S=\{q_1\}$ or $S=\{q_{k-1}\}$, the new vertex extends $Q$ to a $P_k$. Conversely, $Q_S$ has $k$ vertices, so an induced $P_k$ must use all of them. The new vertex then has degree one in that path, and its unique neighbor must be an endpoint of $Q$.
\end{proof}

\begin{lemma}[Prime-free graphs are closed under substitution]\label{lem:substitution}
Let $H$ be prime. If $G_0$ and $G_1$ are $H$-free and $G$ is obtained from $G_0$ by substituting $G_1$ for one vertex, then $G$ is $H$-free.
\end{lemma}

\begin{proof}
An induced copy of $H$ in $G$ either lies inside one substituted part, avoids it, meets it in one vertex, or meets it in at least two vertices and also outside. The first three cases give a copy in $G_1$ or $G_0$. In the last case, the vertices from the substituted part form a nontrivial homogeneous set in the copy, contradicting primeness.
\end{proof}

\begin{definition}
A family $\Pi\subseteq 2^{V(Q)}$ is $P_k$-safe over $Q$ if for every $P_k$-free graph $B$ and every map $\nu:V(B)\to\Pi$, the graph obtained from the disjoint union of $Q$ and $B$ by joining each $b$ to exactly the vertices of $\nu(b)$ is $P_k$-free.
\end{definition}

\begin{theorem}[Exact safe profile classification]\label{thm:safe-profiles}
For $k\ge 4$, a family $\Pi\subseteq 2^{V(Q)}$ is $P_k$-safe over $Q$ if and only if either $\Pi=\emptyset$, or $\Pi=\{S\}$ is a singleton and $S\notin\{\{q_1\},\{q_{k-1}\}\}$.
\end{theorem}

\begin{proof}
If $\Pi$ contains distinct $S,T$, choose $q_i\in S\triangle T$, say $q_i\in S\setminus T$. Let $B=x_1-\cdots-x_{k-1}$, assign profile $S$ to $x_1$, and profile $T$ to $x_2,\dots,x_{k-1}$. Then
\[
  q_i-x_1-x_2-\cdots-x_{k-1}
\]
is an induced $P_k$, a contradiction. Thus a nonempty safe family is a singleton. The forbidden endpoint singletons are excluded by \Cref{lem:one-vertex-extension}.

Conversely, if $\Pi=\{S\}$ and $S$ is not an endpoint singleton, then $Q_S$ is $P_k$-free by \Cref{lem:one-vertex-extension}. Since $P_k$ is prime for $k\ge 4$, substituting any $P_k$-free graph for the new vertex preserves $P_k$-freeness by \Cref{lem:substitution}.
\end{proof}

\subsection{The first genuine splitter}

For a profile $S\subseteq V(Q)$ in an ambient graph, write
\[
  C_S=\{v\notin V(Q):N(v)\cap V(Q)=S\}.
\]
The profile $\{q_{k-2},q_{k-1}\}$ is the terminal safe bag from \Cref{prop:terminal-blowup}. The next profile to the left is rigid, but one more step already supports arbitrary $P_k$-free bipartite interfaces.

\begin{proposition}[Depth-three terminal boundary is rigid]\label{prop:terminal-rigid}
Fix $k\ge 5$ and let $Q=q_1-\cdots-q_{k-1}$ be an induced path in a $P_k$-free graph. Let
\[
  A=C_{\{q_{k-2},q_{k-1}\}},\qquad B=C_{\{q_{k-3},q_{k-2},q_{k-1}\}}.
\]
Then every vertex of $B$ is complete to every vertex of $A$.
\end{proposition}

\begin{proof}
If $b\in B$ and $a\in A$ are nonadjacent, then
\[
  q_1-q_2-\cdots-q_{k-4}-q_{k-3}-b-q_{k-1}-a
\]
is an induced path on $k$ vertices. The profile definitions rule out all possible chords.
\end{proof}

\begin{theorem}[Exact first-splitter realization]\label{thm:first-splitter}
Fix $k\ge 6$ and $Q=q_1-\cdots-q_{k-1}$. Put
\[
  S=\{q_{k-2},q_{k-1}\},\qquad T=\{q_{k-4},q_{k-2},q_{k-1}\}.
\]
For every bipartite graph $R=(A,B;E_R)$, let $G_k(R)$ be obtained from the disjoint union of $Q$ and $A\cup B$ by making $A$ and $B$ stable, adding exactly the interface edges $E_R$ between $A$ and $B$, and assigning profile $S$ to every vertex of $A$ and profile $T$ to every vertex of $B$. Then
\[
  G_k(R)\text{ is }P_k\text{-free}\quad\text{if and only if}\quad R\text{ is }P_k\text{-free}.
\]
Thus $C_S=A$, $C_T=B$, and the interface $G_k(R)[A,B]$ is exactly $R$.
\end{theorem}

\begin{proof}
It is enough to show that no induced $P_k$ can be mixed between $Q$ and $A\cup B$. Write
\[
  p=q_{k-4},\quad r=q_{k-3},\quad s=q_{k-2},\quad t=q_{k-1}.
\]
Every outside vertex sees $s,t$, and vertices of $B$ also see $p$.

Let $P$ be an induced path meeting both $Q$ and $A\cup B$. If $t\in P$, then $s\notin P$, since $s,t$ and any outside vertex form a triangle. The vertex $t$ is adjacent to every outside vertex in $P$, so at most two outside vertices can occur; any attached $Q$ part has at most $q_1,\dots,q_{k-4}$, giving at most $k-1$ vertices. If $t\notin P$ but $s\in P$, then all outside vertices in $P$ see $s$. If $r\in P$, at most one outside vertex can occur, giving at most $k-1$ vertices. If $r\notin P$, the $Q$ part away from $s$ can attach only through $p$ and a vertex of $B$, and again the total is at most $k-1$. Finally, if neither $s$ nor $t$ lies in $P$, the only attachment from $Q$ to the outside is through $p$ and vertices of $B$; the outside tail attached there has length at most two, so the path has at most $k-2$ vertices.

Therefore every induced $P_k$ in $G_k(R)$ lies wholly inside $A\cup B$, where the induced graph is $R$. The converse is immediate because $R$ is an induced subgraph.
\end{proof}

\begin{proposition}[Chain graphs force unbounded first-splitter quotients]\label{prop:chain-quotients}
For every $m\ge 1$ and every $k\ge 6$, the first splitter realizes a $P_k$-free graph whose interface has no side-respecting homogeneous quotient with fewer than $m$ classes on one side.
\end{proposition}

\begin{proof}
Let $F_m$ be the chain graph with bipartition $A=\{a_1,\dots,a_m\}$ and $B=\{b_1,\dots,b_m\}$, with $a_i b_j$ an edge exactly when $j\le i$. This graph is $2K_2$-free, hence $P_6$-free and therefore $P_k$-free for every $k\ge 6$. By \Cref{thm:first-splitter}, it appears as a first-splitter interface.

If two vertices $a_i,a_j$ with $i<j$ lie in the same quotient class on the $A$ side, then $b_j$ distinguishes them. Hence the pair between that $A$-class and the $B$-class containing $b_j$ is mixed, contrary to a homogeneous quotient. Thus $m$ classes are needed on the $A$ side.
\end{proof}

\begin{remark}
Known results on colored $P_6$-free bipartite graphs give large homogeneous bipartite subgraphs in that world \cite{AlecuAtminasLozinZamaraev2021}. \Cref{prop:chain-quotients} is therefore not saying that no pure pair exists. It says that pure pairs alone do not give a bounded finite quotient at the first splitter.
\end{remark}

\section{Cutpoint canonization and the semantic gap}

This section packages a standard-looking ordered canonization lemma and the obstruction that remains after applying it. The lemma itself is not claimed to be new; it is an Erd\H{o}s-Szekeres plus one-switch extraction argument.

\begin{definition}
Let $h,d,t\ge 1$ and let $\Sigma$ be finite. A coloring $\chi$ of ordered $h$-tuples of distinct elements of a finite set $X$ is $(d,t,\Sigma)$-cutpoint-coded if each $x\in X$ has a unary label $\lambda(x)\in\Sigma$ and cutpoints $c_1(x),\dots,c_d(x)\in\{0,1,\dots,t\}$, and $\chi(x_1,\dots,x_h)$ is determined by the labels and the comparison bits
\[
  [c_p(x_i)\le c_q(x_j)]
  \qquad (1\le p,q\le d,
  \ 1\le i<j\le h).
\]
\end{definition}

\begin{lemma}[One-switch extraction]\label{lem:one-switch}
Let $x_1<\cdots<x_N$ be ordered, and let $R$ be a binary relation on pairs $(x_i,x_j)$ with $i<j$. If for each fixed $i$ the set of $j>i$ such that $R(x_i,x_j)$ holds is a suffix of $\{i+1,\dots,N\}$, then there is $I\subseteq\{1,\dots,N\}$ with $|I|\ge \lfloor\sqrt N\rfloor$ such that $R$ is constant on all ordered pairs from $I$.
\end{lemma}

\begin{proof}
Let $m=\lfloor\sqrt N\rfloor$ and split the order into $m$ consecutive blocks of size at least $m$. If some block has no true pair, use it. Otherwise, in each block choose an index $i_r$ whose row has a true entry later in the same block. By the suffix property, $R(x_{i_r},x_{i_s})$ holds for all $r<s$.
\end{proof}

\begin{theorem}[Polynomial ordered canonization]\label{thm:cutpoint-canonization}
For every $d\ge 1$ and finite $\Sigma$ there is $c=c(d,\Sigma)>0$ such that the following holds for all $h,t\ge 1$. If $\chi$ is a $(d,t,\Sigma)$-cutpoint-coded coloring on $X$, then there is an ordered subset
\[
  Y=\{y_1\prec\cdots\prec y_m\}\subseteq X,
  \qquad |Y|\ge c|X|^{2^{-(d^2+d-1)}},
\]
such that all vertices of $Y$ have the same unary label and, for every $p,q$, the bit $[c_p(y_i)\le c_q(y_j)]$ is independent of $i<j$. Consequently $\chi$ is constant on all increasing ordered $h$-tuples from $Y$.
\end{theorem}

\begin{proof}
Pass to a largest unary-label class. Order it by nondecreasing $c_1$. Repeatedly applying the Erd\H{o}s-Szekeres theorem \cite{ErdosSzekeres1935} to $c_2,\dots,c_d$, obtain a subset $X_0$ of size at least a constant times $|X|^{2^{-(d-1)}}$ on which every $c_q$ is monotone.

For each ordered pair $(p,q)$, consider the relation $R_{pq}(x_i,x_j)$ defined by $c_p(x_i)\le c_q(x_j)$. If $c_q$ is nondecreasing along $X_0$, then each row of $R_{pq}$ is a suffix; if $c_q$ is nonincreasing, the complement has this property. Applying \Cref{lem:one-switch} to each of the $d^2$ relations loses a square root each time. The resulting set has the claimed polynomial size and fixes every comparison bit.
\end{proof}

Suppose now that $X=\{x_1<\cdots<x_m\}$ is an ordered block on which such tuple data have canonized. The actual graph edge relation on $X$ may still be an independent binary relation unless the coding is known to decode it.

For $1\le i<j\le h$, define the actual pair-shadow
\[
  S_{ij}(X)=\{1_{a_i a_j\in E(G)}:a_1<\cdots<a_h\text{ in }X\}\subseteq\{0,1\}.
\]

\begin{proposition}[Determinate pair-shadow gives a pure block]\label{prop:determinate-pair-shadow}
If $S_{ij}(X)=\{e\}$ for some $i<j$ and $e\in\{0,1\}$, then $X$ contains $Y$ with $|Y|\ge (|X|-O_h(1))/h$ such that $G[Y]$ is complete when $e=1$ and empty when $e=0$.
\end{proposition}

\begin{proof}
Let $r=i-1$, $g=j-i-1$, and $s=h-j$. Delete the first $r$ and last $s$ vertices of $X$, then choose a maximal subsequence whose consecutive selected vertices have at least $g$ unused vertices between them. Any two selected vertices can be completed to an increasing $h$-tuple in positions $i,j$, so their adjacency is always $e$.
\end{proof}

For an ordered set $X$ and $y\in X$, let $\operatorname{alt}_X(y)$ be the maximum $\ell$ for which there are
\[
  z_0<z_1<\cdots<z_\ell\quad\text{in }X\setminus\{y\}
\]
with $1_{yz_t\in E(G)}\ne 1_{yz_{t+1}\in E(G)}$ for all $t<\ell$.

\begin{lemma}[Bounded alternation is finite cutpoint control]\label{lem:bounded-alt}
If $\max_{y\in X}\operatorname{alt}_X(y)\le A$, then each neighborhood word $z\mapsto 1_{yz\in E(G)}$ is determined by at most $A$ change-points and an initial bit. Conversely, for a fixed $y$, the least number of change-points needed is the number of maximal constant runs minus one.
\end{lemma}

\begin{proof}
The maximal constant runs of the binary word are intervals. An alternating subsequence can take at most one entry from each run, and choosing one from each run gives alternation length equal to the number of runs minus one.
\end{proof}

\begin{proposition}[A local \texorpdfstring{$C_5$}{C5} semantic obstruction]\label{prop:C5-semantic}
For every $m\ge 2$ there is a $C_5$-free graph $G_m$ with a vertex $x$ whose neighborhood alternates $2m-1$ times along a natural order, while both side-patterns needed to extend $x$ to a $C_5$ are present. The obstruction is only the cross matrix between the neighbor and non-neighbor sides.
\end{proposition}

\begin{proof}
Let $V(G_m)=\{x\}\cup P\cup Q$, where $P=\{p_1,\dots,p_m\}$ and $Q=\{q_1,\dots,q_m\}$. Make $P$ independent, $Q$ complete, $P$ complete to $Q$, $x$ complete to $P$, and $x$ anticomplete to $Q$. In the order
\[
  p_1<q_1<p_2<q_2<\cdots<p_m<q_m<x,
\]
the neighborhood word of $x$ alternates.

The graph is $C_5$-free. Any induced cycle avoiding $x$ lies in the split graph on $P\cup Q$, which has no induced $C_5$. If an induced $C_5$ uses $x$, its two neighbors on the cycle lie in $P$. The remaining two vertices cannot both lie in $P$, and any vertex of $Q$ is adjacent to both chosen vertices of $P$, creating a chord or a wrong degree. Still, $P$ contains nonedges and $Q$ contains edges, so the two side-patterns exist. What is missing is a mixed cross matrix, since $P$ is complete to $Q$.
\end{proof}

\begin{remark}
After cutpoint canonization, one rigorous descent route is a determinate pair-shadow, and another is bounded actual-neighborhood alternation. \Cref{prop:C5-semantic} shows that unbounded alternation plus local side patterns does not by itself force the forbidden graph. Any successful descent must use additional $H$-dependent structure.
\end{remark}

\section{Witness blockades: correcting the split condition}

Fix a prime graph $H$ on ordered vertex set $[h]$. For each $i\ge 3$, choose indices $a_i,b_i<i$ such that $a_i i\in E(H)$ and $b_i i\notin E(H)$.

\begin{definition}
A blockade in $G$ is a tuple $\mathcal B=(B_1,\dots,B_h)$ of pairwise disjoint nonempty vertex sets. Its width is $\operatorname{wd}(\mathcal B)=\min_i |B_i|$. A contraction is a tuple $X=(X_1,\dots,X_h)$ with $X_i\subseteq B_i$.

Let $\rho,\lambda>0$ and let $W=\operatorname{wd}(\mathcal B)$. A contraction is $\rho$-admissible if $\min_i |X_i|\ge W^\rho$.
\begin{enumerate}[label=(\alph*)]
\item $\mathcal B$ is $(\rho,\lambda)$-split if for every $\rho$-admissible $X$ and every $i\ge 3$, there is $v\in X_i$ such that
\[
  |N(v)\cap X_{a_i}|\ge m^\lambda,
  \qquad |X_{b_i}\setminus N(v)|\ge m^\lambda,
\]
where $m=\min_j |X_j|$.
\item $\mathcal B$ is $(\rho,\lambda)$-rich if for every $\rho$-admissible $X$ and every $i$, there is $v\in X_i$ such that for every $j<i$,
\[
  \begin{cases}
  |N(v)\cap X_j|\ge m^\lambda, & ij\in E(H),\\
  |X_j\setminus N(v)|\ge m^\lambda, & ij\notin E(H).
  \end{cases}
\]
\end{enumerate}
\end{definition}

\begin{proposition}[Rich blockades force a transversal copy]\label{prop:rich-transversal}
Let $\mathcal B$ be a $(\rho,\lambda)$-rich blockade of width $W$, with $0<\lambda\le 1$ and $\rho\le \lambda^{h-2}$. Then $G$ contains an induced copy of $H$ with one vertex in each block $B_i$.
\end{proposition}

\begin{proof}
Choose vertices backwards. At stage $i$, maintain candidate sets in blocks $1,\dots,i$ of size at least $W^{\lambda^{h-i}}$. Since $\lambda^{h-i}\ge\rho$, richness gives $v_i\in X_i$ with the correct relation to every earlier candidate set at scale $m^\lambda$, leaving sets of size at least $W^{\lambda^{h-i+1}}$. After reaching $i=2$, choose any remaining $v_1$. The construction enforces all adjacencies of $H$.
\end{proof}

The original two-anchor split condition is too weak.

\begin{proposition}[Two-anchor split does not imply rich]\label{prop:two-anchor-not-rich}
Assume $h\ge 4$. For every $W\ge 1$ there are an $H$-free graph $G$ and a width-$W$ blockade $\mathcal B=(B_1,\dots,B_h)$ such that $\mathcal B$ is $(\rho,\lambda)$-split for every $0<\rho,\lambda\le 1$, but no positive-width contraction is rich for any parameters.
\end{proposition}

\begin{proof}
Choose $j_0\in\{1,2,3\}\setminus\{a_4,b_4\}$. Let $Q$ be obtained from $H$ by toggling only the relation between $4$ and $j_0$. Replace each vertex $i$ of $Q$ by an independent set $B_i$ of size $W$, making block pairs complete or anticomplete according to $Q$.

The chosen witness relations $i a_i$ and $i b_i$ are unchanged for every $i\ge 3$, so the blockade is split at all positive scales. But every vertex of $B_4$ has the wrong relation to $B_{j_0}$, so no positive-width contraction can be rich.

The graph is $H$-free. Any induced copy of prime $H$ in this blow-up is either contained in one block, impossible by size, or is transversal. Every transversal $h$-set induces $Q$, not $H$.
\end{proof}

The replacement target is coordinatewise split.

\begin{definition}
For a contraction $X$ and $v\in X_i$, define the correct-coordinate degree into $X_j$ by
\[
 d^H_j(v;X)=
 \begin{cases}
 |N(v)\cap X_j|, & ij\in E(H),\\
 |X_j\setminus N(v)|, & ij\notin E(H).
 \end{cases}
\]
The blockade is coordinatewise $(\rho,\lambda)$-split if for every $\rho$-admissible contraction $X$ of width $m$, every $i$ and every $j<i$, some $v\in X_i$ satisfies $d^H_j(v;X)\ge m^\lambda$.
\end{definition}

\begin{lemma}[A failed coordinate cleans to a wrong pure pair]\label{lem:failed-coordinate}
Let $0<\lambda<1$ and let $X$ be a contraction of width $m$. Fix $i>j$. If every $v\in X_i$ satisfies $d^H_j(v;X)<m^\lambda$, then for every $0<\theta<1-\lambda$ and all sufficiently large $m$, there are $X'\subseteq X_i$ and $Y'\subseteq X_j$ with $|X'|,|Y'|\ge m^\theta$ such that every pair between $X'$ and $Y'$ has the wrong relation relative to $H$.
\end{lemma}

\begin{proof}
Choose $X'\subseteq X_i$ of size $\lfloor m^\theta\rfloor$. Let $U$ be the union, over $x\in X'$, of the vertices of $X_j$ having the correct relation to $x$. Then $|U|<m^{\theta+\lambda}=o(m)$. Choose $Y'\subseteq X_j\setminus U$ of size $\lfloor m^\theta\rfloor$.
\end{proof}

\begin{proposition}[Coordinatewise split synchronizes to rich]\label{prop:coordinate-synchronizes}
Fix $0<\rho_0<\rho_1\le 1$ and $0<\lambda'<\lambda\le 1$. For all sufficiently large widths, every coordinatewise $(\rho_0,\lambda)$-split blockade is $(\rho_1,\lambda')$-rich.
\end{proposition}

\begin{proof}
Let $X$ be a $\rho_1$-admissible contraction of width $m$, and put $s=\lfloor m/h\rfloor$. For large width, $s\ge W^{\rho_0}$ and $s^\lambda\ge m^{\lambda'}$.

For fixed $i$ and $j<i$, let
\[
  T_{ij}(s)=\{v\in X_i:d^H_j(v;X)<s^\lambda\}.
\]
If $|T_{ij}(s)|\ge s$, then taking $s$ vertices from $T_{ij}(s)$ and $s$ vertices from each other block gives a $\rho_0$-admissible contraction contradicting coordinatewise split. Thus $|T_{ij}(s)|<s$. Since the union of the bad sets over $j<i$ has size less than $(i-1)s\le (h-1)s<m$, some $v\in X_i$ avoids all of them, and then $d^H_j(v;X)\ge s^\lambda\ge m^{\lambda'}$ for every $j<i$.
\end{proof}

\begin{remark}
The missing theorem is not ``two-anchor split implies rich.'' That is false by \Cref{prop:two-anchor-not-rich}. A viable blockade route must extract coordinatewise split, or turn failed coordinates into useful wrong pure pairs or quotient obstructions.
\end{remark}

\section{The rooted \texorpdfstring{$P_4$}{P4} two-anchor branch}

This final section isolates a local configuration that appears when trying to extend a rooted $P_4$ to a $C_5$-type obstruction. The goal is to separate the part that is elementary two-anchor extraction from the part that genuinely needs $H$-dependent compression.

\begin{definition}[Purified rooted \texorpdfstring{$P_4$}{P4} configuration]\label{def:purified}
A purified rooted $P_4$ configuration in a graph $G$ is a tuple of pairwise disjoint sets $(B_1,B_2,B_3,B_4,Y)$ such that
\[
  B_1\text{ is anticomplete to }B_3\cup B_4,
  \quad B_2\text{ is anticomplete to }B_4,
  \quad Y\text{ is anticomplete to }B_2\cup B_3.
\]
For an edge $e=b_2b_3\in E(B_2,B_3)$, define
\[
  A_e=N(b_2)\cap B_1,
  \qquad D_e=N(b_3)\cap B_4,
  \qquad w(e)=\min\{|A_e|,|D_e|\}.
\]
Let $T=\sum_{e\in E(B_2,B_3)}w(e)$ and $\Lambda=|E(B_2,B_3)|$.
\end{definition}

\begin{proposition}[Two-anchor extraction with a defect term]\label{prop:defect-extraction}
Let $(B_1,B_2,B_3,B_4,Y)$ be purified. Fix $e=b_2b_3\in E(B_2,B_3)$ and subsets $A\subseteq A_e$, $D\subseteq D_e$ with $|A|=|D|=m$. Let
\[
  Y^*(A,D;e)=\{y\in Y:N(y)\cap A\ne\emptyset\ne N(y)\cap D\}.
\]
Then $G$ contains an anticomplete pair $(X,Z)$ with
\[
  |X|\ge m,
  \qquad |Z|\ge (|Y|-|Y^*(A,D;e)|)/2.
\]
\end{proposition}

\begin{proof}
Let $Y_A=\{y\in Y:N(y)\cap A=\emptyset\}$ and $Y_D=\{y\in Y:N(y)\cap D=\emptyset\}$. Since $Y\setminus Y^*\subseteq Y_A\cup Y_D$, one of $Y_A,Y_D$ has size at least $(|Y|-|Y^*|)/2$. Pair it with $A$ or $D$.
\end{proof}

\begin{theorem}[Purified two-anchor extraction]\label{thm:two-anchor-extraction}
Let $(B_1,B_2,B_3,B_4,Y)$ be purified. Assume there is no induced cycle $a-b_2-b_3-d-y-a$ with $(a,b_2,b_3,d,y)\in B_1\times B_2\times B_3\times B_4\times Y$. If $\Lambda>0$, then $G$ contains an anticomplete pair $(X,Z)$ with
\[
  |X|\ge T/\Lambda,
  \qquad |Z|\ge |Y|/2.
\]
If $\Lambda=0$, then $(B_2,B_3)$ is already anticomplete.
\end{theorem}

\begin{proof}
If $\Lambda=0$, there is nothing to prove. Otherwise choose $e=b_2b_3$ with $w(e)\ge T/\Lambda$, and choose $A\subseteq A_e$, $D\subseteq D_e$ of size $w(e)$. If $y\in Y^*(A,D;e)$, then some $a\in A$ and $d\in D$ give the induced cycle $a-b_2-b_3-d-y-a$, using the purification conditions. Hence $Y^*=\emptyset$, and \Cref{prop:defect-extraction} applies.
\end{proof}

\begin{proposition}[Trace-count compression forces either a pair or large supply]\label{prop:trace-count}
Let $0<\eta<\beta$ and let $(B_1,B_2,B_3,B_4,Y)$ be purified with $|B_i|\ge n^\beta$ for $1\le i\le 4$. Define endpoint trace families
\[
  \mathcal T_1=\{N(b_2)\cap B_1:b_2\in B_2\},
  \qquad
  \mathcal T_4=\{N(b_3)\cap B_4:b_3\in B_3\}.
\]
If $|\mathcal T_1|,|\mathcal T_4|\le n^\eta$, then either:
\begin{enumerate}[label=(\roman*)]
\item $G$ contains an anticomplete pair $(X,Z)$ with $|X|,|Z|\ge n^{\beta-\eta}$; or
\item there is an edge $b_2b_3\in E(B_2,B_3)$ with
\[
  |N(b_2)\cap B_1|\ge |B_1|-n^{\beta-\eta},
  \qquad
  |N(b_3)\cap B_4|\ge |B_4|-n^{\beta-\eta}.
\]
\end{enumerate}
\end{proposition}

\begin{proof}
Choose a largest fiber $U\subseteq B_2$ for the trace into $B_1$, and a largest fiber $V\subseteq B_3$ for the trace into $B_4$. Then $|U|,|V|\ge n^{\beta-\eta}$. If $U$ and $V$ are anticomplete, (i) holds. Otherwise choose an edge $b_2b_3$ between them. If the common trace $A=N(b_2)\cap B_1$ misses at least $n^{\beta-\eta}$ vertices of $B_1$, then $(U,B_1\setminus A)$ is anticomplete. Similarly on the right. If neither anticomplete pair appears, (ii) holds.
\end{proof}

\subsection{Crossing traces and rectangle shattering}

For $y\in Y$, define
\[
  P_y=\{b_2\in B_2: \exists a\in B_1\cap N(y)\cap N(b_2)\},
  \qquad
  Q_y=\{b_3\in B_3: \exists d\in B_4\cap N(y)\cap N(b_3)\},
\]
and
\[
  C_y=E(B_2,B_3)\cap (P_y\times Q_y).
\]
Thus $C_y$ is the set of middle edges that can participate in the cycle $a-b_2-b_3-d-y-a$ with this $y$.

\begin{theorem}[Star forests are exactly the shattered rectangle patterns]\label{thm:star-forest}
Let $M$ be a bipartite graph with bipartition $(U,V)$, and let
\[
  \mathcal R(M)=\{E(M)\cap(P\times Q):P\subseteq U,
  \ Q\subseteq V\}.
\]
For a finite $F\subseteq E(M)$, the following are equivalent:
\begin{enumerate}[label=(\roman*)]
\item $F$ is shattered by $\mathcal R(M)$;
\item the graph $(V(F),F)$ is a star forest.
\end{enumerate}
\end{theorem}

\begin{proof}
If $F$ contains a three-edge path $u_1-v_1-u_2-v_2$, then the subset $\{u_1v_1,u_2v_2\}$ cannot be realized by a rectangle: including those two edges forces $u_1,u_2$ and $v_1,v_2$ to be selected, hence also includes $u_2v_1$. Thus every component of $F$ has no three-edge path, so it is a star.

Conversely, if $F$ is a star forest, realize any subset of $F$ component by component. For a star centered in $U$, put the center in $P$ if some chosen edge of that component is desired, and put exactly the desired leaves in $Q$; do the symmetric construction for stars centered in $V$.
\end{proof}

\begin{corollary}[Star-forest exclusion compresses crossing traces]\label{cor:star-compress}
Let $\mathcal C=\{C_y:y\in Y\}\subseteq 2^{E(B_2,B_3)}$. If $\mathcal C$ shatters no star forest with $t$ edges inside $G[B_2,B_3]$, then
\[
  |\mathcal C|\le \sum_{i=0}^{t-1}\binom{\Lambda}{i}=O_t(\Lambda^{t-1}).
\]
\end{corollary}

\begin{proof}
Any set shattered by $\mathcal C$ is shattered by the full rectangle system, hence is a star forest by \Cref{thm:star-forest}. Therefore $\VCdim(\mathcal C)<t$, and the Sauer-Shelah lemma applies \cite{Sauer1972,Shelah1972}.
\end{proof}

\begin{proposition}[Boolean local obstructions]\label{prop:boolean-obstructions}
For every $m\ge 1$ there are purified rooted $P_4$ configurations with the following properties.
\begin{enumerate}[label=(\roman*)]
\item The crossing trace family shatters a star with $m$ edges, while $T=\Lambda=m$ and hence $T/\Lambda=1$.
\item The crossing trace family shatters a matching with $m$ edges, and both one-sided trace families $\{P_y:y\in Y\}$ and $\{Q_y:y\in Y\}$ have size $2^m$.
\end{enumerate}
\end{proposition}

\begin{proof}
For (i), take $B_1=\{a\}$, $B_2=\{u\}$, $B_3=\{v_1,\dots,v_m\}$, and $B_4=\{d_1,\dots,d_m\}$. Add middle edges $uv_i$ and endpoint edges $au$, $v_i d_i$. For each $S\subseteq[m]$, add $y_S$ adjacent to $a$ and to $d_i$ exactly for $i\in S$. Then $C_{y_S}=\{uv_i:i\in S\}$ and every middle edge has weight one.

For (ii), take $B_1=\{a_1,\dots,a_m\}$, $B_2=\{u_1,\dots,u_m\}$, $B_3=\{v_1,\dots,v_m\}$, and $B_4=\{d_1,\dots,d_m\}$. Add middle edges $u_iv_i$ and endpoint edges $a_iu_i$, $v_id_i$. For each $S\subseteq[m]$, add $y_S$ adjacent to $a_i$ and $d_i$ exactly for $i\in S$. Then $P_{y_S}=\{u_i:i\in S\}$, $Q_{y_S}=\{v_i:i\in S\}$, and $C_{y_S}=\{u_iv_i:i\in S\}$.
\end{proof}

The previous examples are local. The next statement makes the obstruction $H$-specific through two explicit hereditary template closures.

\begin{definition}[Crossing-cube template closures]\label{def:cube-closures}
Let $\mathcal P$ be the hereditary class of graphs admitting a partition
\[
  X\cup L\cup R\cup D\cup Y,
  \qquad |D|\le 1,
\]
such that $X,L,R,Y$ are stable, the only allowed edges are between $X-L$, $L-R$, $R-D$, $D-Y$, and $X-Y$, the graph between $L$ and $R$ is a matching, the graph between $X$ and $Y$ is a matching, the unique vertex of $D$, if present, is complete to $R\cup Y$, and the graph between $X$ and $L$ is arbitrary.

Let $\mathcal B$ be the hereditary class of graphs admitting a partition
\[
  A\cup C\cup R\cup Z\cup Y,
  \qquad |A|\le 1,
  \quad |C|\le 1,
\]
such that $R,Z,Y$ are stable, the only allowed edges are between $A-C$, $C-R$, $R-Z$, $Z-Y$, and $A-Y$, the edge $AC$ is present when both singletons exist, $C$ is complete to $R$, $A$ is complete to $Y$, the graph between $Z$ and $Y$ is a matching, and the graph between $R$ and $Z$ is arbitrary.
\end{definition}

\begin{theorem}[H-specific crossing-cube obstruction]\label{thm:H-specific-cubes}
Let $H$ be fixed.
\begin{enumerate}[label=(\roman*)]
\item If $H\notin\mathcal P$, then for every $m$ there is an $H$-free purified rooted $P_4$ configuration with $T/\Lambda=1$ whose crossing trace family shatters a matching of size $m$.
\item If $H\notin\mathcal B$, then for every $m$ there is an $H$-free purified rooted $P_4$ configuration with $T/\Lambda=1$ whose crossing trace family shatters a star of size $m$.
\end{enumerate}
Consequently, an $H$-free theorem forcing bounded star-forest shattering in all purified rooted $P_4$ configurations cannot follow from the rooted $P_4$ axioms alone unless $H\in\mathcal P\cap\mathcal B$.
\end{theorem}

\begin{proof}
Assume $H\notin\mathcal P$. Define the private-matching cube $PM_m$ by
\[
  B_1=\{x_S:S\subseteq[m]\},\quad B_2=\{\ell_i:i\in[m]\},\quad B_3=\{r_i:i\in[m]\},\quad B_4=\{d\},\quad Y=\{y_S:S\subseteq[m]\}.
\]
The edges are exactly $x_S\ell_i$ when $i\in S$, $\ell_i r_i$ for all $i$, $r_i d$ for all $i$, $dy_S$ for all $S$, and $y_Sx_S$ for all $S$. This is purified. The middle graph is a matching, every middle edge has weight one, and $C_{y_S}=\{\ell_i r_i:i\in S\}$. Thus $T=\Lambda=m$ and the matching is shattered. The whole graph belongs to $\mathcal P$, so it is $H$-free because $\mathcal P$ is hereditary and $H\notin\mathcal P$.

Assume $H\notin\mathcal B$. Define the Boolean-star cube $BS_m$ by
\[
  B_1=\{a\},\quad B_2=\{c\},\quad B_3=\{r_i:i\in[m]\},\quad B_4=\{z_S:S\subseteq[m]\},\quad Y=\{y_S:S\subseteq[m]\}.
\]
The edges are exactly $ac$, $cr_i$ for all $i$, $r_iz_S$ when $i\in S$, $z_Sy_S$ for all $S$, and $y_Sa$ for all $S$. This is purified. The middle graph is a star, every middle edge has weight one, and $C_{y_S}=\{cr_i:i\in S\}$. The graph belongs to $\mathcal B$, hence is $H$-free under the hypothesis.
\end{proof}

\begin{remark}
Even when $H\in\mathcal P\cap\mathcal B$, more is needed. Arbitrary shattering gives traces, but it need not give the synchronized private witnesses present in $PM_m$ or $BS_m$. This is why an endpoint or crossing compression lemma should be stated with the exact synchronization it uses.
\end{remark}

\section{Questions left by the trimmed note}

The results above leave the following concrete targets. These are deliberately phrased as questions rather than conjectures.

\begin{question}[Affine occupancy with an H-dependent escape]
Can one prove an $H$-dependent level-two shadow theorem saying that either the descended exact rows have bounded affine-cell occupancy, or the sparse-occupancy obstruction already yields a polynomial complete or anticomplete pair? \Cref{thm:singleton-cut} shows that the unconditional occupancy statement is false.
\end{question}

\begin{question}[Path quotient or bypass]
Can one compress the finite colored obstruction data beyond the terminal safe bag into finitely many quotient states, or bypass it directly with a polynomial pure pair? \Cref{thm:first-splitter,prop:chain-quotients} show that the first splitter already realizes nested-neighborhood diversity that cannot be captured by a bounded homogeneous quotient.
\end{question}

\begin{question}[Coordinatewise blockade extraction]
Can a prime $H$-free graph with no polynomial clique or stable set always be refined to a polynomial-width blockade that is coordinatewise split for parameters strong enough to imply richness? By \Cref{prop:coordinate-synchronizes}, this would force a transversal copy of $H$; by \Cref{prop:two-anchor-not-rich}, the older two-anchor split target is insufficient.
\end{question}

\begin{question}[Semantic descent after cutpoint canonization]
After \Cref{thm:cutpoint-canonization}, can the large ordered canonical block always be sent to either a proper induced-subgraph reduction, a determinate pair-shadow, bounded actual-neighborhood alternation, or another finite cutpoint control of the actual edge relation? \Cref{prop:C5-semantic} shows that side patterns plus unbounded alternation are not enough.
\end{question}

\begin{question}[Endpoint and crossing trace compression]
In an $H$-free purified rooted $P_4$ configuration with no polynomial pure pair, can one force endpoint trace compression, bounded star-forest shattering of crossing traces, or a weaker supply condition sufficient for \Cref{thm:two-anchor-extraction}? By \Cref{thm:H-specific-cubes}, any unconditional bounded-shattering theorem must use information about $H$ beyond the rooted $P_4$ axioms.
\end{question}

\section*{Acknowledgement of uncertainty}
The elementary lemmas in this note may well overlap with existing folklore in extremal graph theory, stable regularity, VC theory, or the induced-density literature. The intended contribution, if any, is the explicit packaging of negative examples against particular reductions. The most useful response from readers would be pointers to known versions, stronger known obstructions, or mistakes in the formulations above.

\printbibliography

\end{document}
